\chapter{Fundamentação teórica}
\label{chap:fundamentacao}

Este capítulo apresenta os conceitos em que o trabalho se apoia. A \autoref{sec:microplastico} trata do microplástico e do seu caminho da terra até o mar. A \autoref{sec:bacias} discute o delineamento de bacias hidrográficas a partir do relevo e o algoritmo de percolação invasiva usado aqui. A \autoref{sec:cca} define cidades a partir da densidade populacional. A \autoref{sec:leis-potencia} apresenta as leis de potência, a lei de Zipf e os métodos estatísticos para estimá-las, e a \autoref{sec:escala} trata das leis de escala urbana, que dão a forma da relação testada nesta dissertação.

\section{Plástico e microplástico no ambiente}
\label{sec:microplastico}

\subsection{Origem, fragmentação e definição}

Plásticos são polímeros sintéticos de baixo custo, leves e duráveis. A durabilidade que os torna úteis é também o que os torna um problema ambiental: expostos à radiação ultravioleta, à abrasão e às variações de temperatura, os objetos plásticos se fragmentam em pedaços cada vez menores, mas a mineralização completa do polímero leva de décadas a séculos \cite{barnes2009,andrady2011}. O resultado é um acúmulo contínuo de fragmentos no ambiente. \citeonline{thompson2004} mostraram que fragmentos microscópicos de polímeros comuns já eram abundantes em sedimentos e na coluna d'água das costas do Reino Unido nos anos 1960 e que a sua abundância cresceu nas décadas seguintes.

Não há uma definição única de microplástico, mas o limite superior de 5 mm é o mais usado \cite{frias2019}. Costuma-se distinguir os microplásticos \textbf{secundários}, que resultam da fragmentação de objetos maiores, dos \textbf{primários}, fabricados já nesse tamanho. Entre os primários estão os \textit{pellets} de resina pré-produção, conhecidos em inglês como \textit{nurdles}: grânulos de 2 a 5 mm que são a matéria-prima da indústria de transformação de plásticos. Os \textit{pellets} são perdidos durante a produção, o transporte e o manuseio, e chegam ao mar pelos sistemas de drenagem e pelos rios \cite{karlsson2018}. Além de serem poluentes em si, eles adsorvem compostos orgânicos persistentes da água, concentrando-os em até seis ordens de grandeza \cite{mato2001}.

\subsection{Rios como vetores de plástico para o oceano}

A maior parte do plástico encontrado no mar tem origem terrestre. \citeonline{jambeck2015} combinaram dados de geração de resíduos, população costeira e gestão de resíduos de 192 países para estimar que entre 4,8 e 12,7 milhões de toneladas de plástico entraram no oceano em 2010. Os rios são a principal via desse transporte. \citeonline{lebreton2017} desenvolveram um modelo global de emissão de plástico por rios em que a entrada de plástico em cada bacia é função da população, da geração de resíduos mal geridos e do escoamento superficial; o modelo estimou entre 1,15 e 2,41 milhões de toneladas por ano e concentrou a maior parte das emissões em poucos rios asiáticos. \citeonline{schmidt2017} chegaram a uma conclusão semelhante ao ajustar uma regressão entre a carga de plástico medida em dezenas de rios e a quantidade de resíduos mal geridos nas suas bacias. \citeonline{meijer2021} refinaram esse tipo de modelo com uma representação probabilística do transporte dentro da bacia, que leva em conta a distância até o rio, o uso do solo, a declividade e a precipitação, e concluíram que mais de mil rios respondem por 80\% das emissões globais, com um peso muito maior de rios pequenos e urbanos do que o estimado antes.

Esses modelos têm em comum a hipótese de que \textbf{a população da bacia é uma variável explicativa central} da carga de plástico que chega à foz. Medições de campo que testem essa hipótese diretamente, bacia a bacia e numa escala continental, são escassas. As evidências empíricas mais diretas vêm de estudos locais, discutidos no \autoref{chap:trabalhos-relacionados}.

\subsection{Medição de microplástico e o \textit{Nurdle Patrol}}

Os métodos de amostragem de microplástico variam muito: redes de arrasto na superfície do mar (com resultados em partículas por m$^3$ ou por km$^2$), testemunhos e amostras de sedimento (em partículas por kg de sedimento seco) e coletas em praias (em partículas por área ou por tempo de busca) \cite{hidalgoruz2012}. Resultados obtidos com métodos diferentes não são diretamente comparáveis, e por isso uma análise quantitativa precisa se restringir a um único método.

O \textit{Nurdle Patrol} é um programa de ciência cidadã coordenado pelo Mission-Aransas National Estuarine Research Reserve, da Universidade do Texas, em que voluntários contam os \textit{pellets} que encontram numa praia durante 10 minutos de busca \cite{tunnell2020}. A unidade resultante, \textit{pellets} por 10 minutos, é uma medida padronizada de esforço: o mesmo protocolo é seguido em todos os pontos, o que torna as contagens comparáveis entre si. Em compensação, a cobertura espacial depende de onde há voluntários, o que concentra as coletas na costa do Golfo do México \cite{tunnell2020}. As contagens do programa estão incluídas no banco de microplásticos marinhos do \gls{NCEI}, da NOAA, que reúne amostras de diversos métodos e fontes com as respectivas unidades \cite{nyadjro2023}.

\section{Bacias hidrográficas e o seu delineamento}
\label{sec:bacias}

\subsection{Bacias a partir de modelos digitais de elevação}

Uma bacia hidrográfica é a região cujo escoamento superficial converge para uma mesma saída. Em termos computacionais, o problema de delineamento de bacias consiste em, dado um \gls{MDE} representado como uma grade de altitudes, atribuir cada célula da grade à saída para onde a água que cai nela escoaria.

A abordagem clássica é a das \textbf{direções de fluxo}. No método D8, cada célula drena para a vizinha, entre as oito, com maior declividade descendente \cite{ocallaghan1984}; variações como o D$\infty$ distribuem o fluxo entre duas vizinhas \cite{tarboton1997}. Acumulando o fluxo ao longo dessas direções obtém-se a área de drenagem de cada célula, e a rede de drenagem é definida pelas células cuja área acumulada passa de um limiar. Essa abordagem tem duas dificuldades para o problema desta dissertação. Primeiro, os MDEs reais têm depressões e platôs em que a direção de fluxo não é definida, e é preciso preenchê-los antes, com algoritmos próprios. Segundo, o resultado natural é uma rede de drenagem, não um conjunto de bacias associadas a saídas no oceano, e a escolha do limiar de área crítica é arbitrária.

Redes fluviais têm estrutura estatística própria: a área das bacias, o comprimento dos rios e a sua ramificação seguem leis de potência com expoentes aproximadamente universais \cite{rodrigueziturbe1997,maritan1996}. Essa observação é usada no \autoref{chap:resultados} para verificar se as bacias delineadas têm a distribuição de tamanhos esperada.

\subsection{Percolação invasiva}

A percolação invasiva foi proposta por \citeonline{wilkinson1983} para modelar o deslocamento de um fluido por outro num meio poroso. Cada sítio do meio tem uma resistência aleatória; a partir de um sítio inicial, o fluido invasor avança sempre pelo sítio de \textbf{menor resistência} da fronteira da região já invadida. Aplicada a um relevo, a resistência de um sítio é a sua altitude: a água que parte de um ponto avança pelo ponto mais baixo ao seu redor, enche as depressões que encontra até transbordar pelo ponto mais baixo da borda e continua descendo. É exatamente o comportamento de uma gota que escoa sobre o terreno, incluindo a travessia de lagos e depressões, sem necessidade de preenchimento prévio.

\citeonline{fehr2009} usaram essa ideia para construir um algoritmo de delineamento de bacias, o \acrshort{IPBA}. Dado um conjunto de saídas rotuladas (\textit{sinks}), o algoritmo percorre as células da grade e, para cada célula ainda sem bacia, executa uma percolação invasiva a partir dela. A percolação para quando atinge um \textit{sink} ou uma célula que já pertence a uma bacia; todas as células invadidas nessa execução recebem então o rótulo da bacia atingida. O \autoref{alg:ipba} resume o procedimento. A fronteira da região invadida é mantida num \textit{heap} binário de mínimo ordenado pela altitude, de modo que cada passo de invasão custa $O(\log m)$, em que $m$ é o tamanho da fronteira. Como cada célula é rotulada uma única vez e as percolações param ao encontrar células já rotuladas, o custo total cresce aproximadamente como $O(n \log n)$ para uma grade de $n$ células.

\begin{algorithm}[ht!]
	\SetSpacedAlgorithm
	\caption{\label{alg:ipba}Delineamento de bacias por percolação invasiva (IPBA)}
	\Entrada{grade de altitudes $h$; rótulos $\lambda(s) > 0$ para os \textit{sinks}; vizinhança de 4 células}
	\Saida{rótulo de bacia $b(k)$ para cada célula de terra $k$}
	\Inicio{
		\ParaCada{célula de terra $k$ ainda sem bacia}{
			$H \leftarrow$ \textit{heap} de mínimo com $(h(k), k)$; \quad $I \leftarrow \emptyset$\;
			\Enqto{$H$ não estiver vazio}{
				$t \leftarrow$ célula de menor altitude extraída de $H$; \quad $I \leftarrow I \cup \{t\}$\;
				\Se{$t$ é um \textit{sink} ou $t$ já tem bacia}{
					$w \leftarrow$ bacia de $t$ (ou o próprio rótulo, se \textit{sink})\;
					\textbf{interrompe}\;
				}
				insere em $H$ os vizinhos de $t$ ainda não visitados, com as suas altitudes\;
			}
			\ParaCada{$i \in I$}{$b(i) \leftarrow w$\;}
		}
	}
\end{algorithm}

O IPBA tem duas propriedades úteis para esta dissertação. A primeira é que ele aceita \textbf{qualquer conjunto de saídas}: se todas as células de terra vizinhas ao oceano forem declaradas \textit{sinks}, o algoritmo produz uma partição completa do continente em bacias costeiras, sem que seja preciso saber de antemão onde ficam os rios. As bacias dos rios grandes emergem sozinhas, porque todas as células a montante escoam até a célula da foz. A segunda é que ele trata depressões e platôs sem pré-processamento. \citeonline{fehr2011} usaram o mesmo arcabouço para estudar a geometria dos divisores de água, que são fractais.

\section{Definição de cidades pela densidade: o \textit{City Clustering Algorithm}}
\label{sec:cca}

Os dados oficiais de população são agregados por unidades administrativas, como municípios e setores censitários. As fronteiras dessas unidades foram traçadas por motivos históricos e políticos, e uma região metropolitana contínua pode estar repartida em dezenas de municípios, enquanto um município grande pode conter vasta área rural. Qualquer estatística sobre o tamanho de cidades depende, portanto, de uma escolha prévia do que é uma cidade \cite{rozenfeld2011,arcaute2015}.

O \gls{CCA} foi proposto por \citeonline{rozenfeld2008} como uma definição de cidade construída de baixo para cima, a partir da distribuição espacial da população. Sobre uma grade de células com população conhecida, o algoritmo tem dois parâmetros:

\begin{itemize}
	\item o \textbf{limiar de densidade} $D^{*}$: uma célula é \textit{ativa} (urbana) se a sua densidade populacional for maior que $D^{*}$;
	\item a \textbf{distância de ligação} $\ell$: duas células ativas pertencem à mesma mancha se estiverem a menos de $\ell$ uma da outra, direta ou indiretamente, por uma cadeia de células ativas.
\end{itemize}

Formalmente, seja $V = \{k : \rho_k > D^{*}\}$ o conjunto de células ativas, com $\rho_k$ a densidade da célula $k$, e seja $G_\ell$ o grafo sobre $V$ com uma aresta entre $j$ e $k$ sempre que a distância geodésica $d(j,k)$ entre os centros das células for menor que $\ell$. As manchas urbanas são as componentes conexas de $G_\ell$, e a população e a área de uma mancha são as somas das populações e áreas das suas células. O CCA é, portanto, um problema de percolação contínua: $D^{*}$ controla a ocupação dos sítios e $\ell$ controla o alcance das ligações \cite{rozenfeld2008,rozenfeld2011}.

A escolha de $D^{*}$ e $\ell$ muda o resultado. Com $D^{*}$ baixo e $\ell$ alto, regiões inteiras se fundem numa única mancha; com $D^{*}$ alto e $\ell$ baixo, as cidades se fragmentam em bairros. \citeonline{rozenfeld2011} mostraram que, para uma ampla faixa de parâmetros, as manchas dos Estados Unidos e do Reino Unido seguem a lei de Zipf, mas que o expoente medido depende dos valores escolhidos. \citeonline{oliveira2014} usaram o CCA para estudar a relação entre o tamanho das cidades e as suas emissões de CO$_2$. Nesta dissertação, $D^{*}$ e $\ell$ são escolhidos pela própria lei de Zipf, como descrito no \autoref{chap:metodologia}.

\section{Leis de potência e a lei de Zipf}
\label{sec:leis-potencia}

\subsection{Leis de potência}

Uma variável aleatória contínua $x$ segue uma lei de potência, a partir de um valor mínimo $x_{\min}$, quando a sua densidade de probabilidade é
\begin{equation}
	p(x) = \frac{\gamma - 1}{x_{\min}} \left( \frac{x}{x_{\min}} \right)^{-\gamma}, \qquad x \geq x_{\min},
	\label{eq:lei-potencia}
\end{equation}
com $\gamma > 1$. A função de distribuição complementar é $P(X \geq x) = (x / x_{\min})^{-(\gamma - 1)}$, uma reta de inclinação $-(\gamma-1)$ num gráfico log-log. Leis de potência aparecem no tamanho de cidades, na frequência de palavras, na intensidade de terremotos e na área de bacias hidrográficas, entre muitos outros fenômenos \cite{newman2005}. Elas não têm escala característica: a razão entre a frequência de eventos de tamanho $2x$ e de tamanho $x$ é a mesma para qualquer $x$.

\subsection{A lei de Zipf}

A lei de Zipf afirma que, ordenando as cidades de um país da maior para a menor, a população $S$ da cidade de posto $r$ segue
\begin{equation}
	S(r) \propto r^{-\alpha}, \qquad \alpha \approx 1,
	\label{eq:zipf}
\end{equation}
ou seja, a segunda maior cidade tem cerca de metade da população da maior, a terceira cerca de um terço, e assim por diante \cite{zipf1949}. A forma \textit{rank-size} da \autoref{eq:zipf} é equivalente a uma lei de potência na distribuição dos tamanhos. Numa amostra de $n$ cidades, o posto da cidade de tamanho $S$ é aproximadamente $r \approx n\, P(X \geq S) \propto S^{-(\gamma - 1)}$; invertendo, $S \propto r^{-1/(\gamma - 1)}$, e portanto
\begin{equation}
	\alpha = \frac{1}{\gamma - 1}, \qquad \text{de modo que} \qquad \alpha = 1 \iff \gamma = 2 .
	\label{eq:alpha-gamma}
\end{equation}

\citeonline{gabaix1999} mostrou que a lei de Zipf emerge se as cidades crescem de forma proporcional, com taxas de crescimento independentes do tamanho (lei de Gibrat), e há uma barreira refletora para as cidades pequenas. A validade da lei, porém, é debatida. \citeonline{soo2005} testou 73 países e rejeitou $\alpha = 1$ na maioria deles. \citeonline{eeckhout2004} mostrou que, quando se usam todos os lugares povoados dos Estados Unidos e não apenas as maiores cidades, a distribuição é mais bem descrita por uma lognormal, e que a lei de Zipf vale apenas para a cauda superior. \citeonline{jiang2011} encontraram a lei de Zipf nas ``cidades naturais'' dos Estados Unidos, definidas pela densidade de cruzamentos de ruas, o que reforça que a definição de cidade importa tanto quanto os dados.

\subsection{Estimação do expoente}
\label{sec:estimacao}

Ajustar uma reta ao gráfico log-log por mínimos quadrados é a forma mais comum de estimar expoentes de leis de potência, mas é um estimador enviesado e os seus erros-padrão não são válidos, porque os pontos do gráfico não são independentes \cite{clauset2009}. Dois estimadores são usados nesta dissertação.

\textbf{Método de Clauset, Shalizi e Newman.} \citeonline{clauset2009} propuseram um procedimento em três passos, que se tornou o padrão para dados empíricos. Primeiro, para um $x_{\min}$ fixo, o expoente é estimado por máxima verossimilhança,
\begin{equation}
	\hat{\gamma} = 1 + n \left[ \sum_{i=1}^{n} \ln \frac{x_i}{x_{\min}} \right]^{-1},
	\qquad \sigma_{\hat{\gamma}} = \frac{\hat{\gamma} - 1}{\sqrt{n}},
	\label{eq:mle}
\end{equation}
em que a soma percorre as $n$ observações com $x_i \geq x_{\min}$. Segundo, $x_{\min}$ é escolhido como o valor que minimiza a distância de \gls{KS} entre a distribuição acumulada empírica da cauda e a do modelo ajustado. Terceiro, a plausibilidade da lei de potência é avaliada por \textit{bootstrap} semiparamétrico: geram-se amostras sintéticas do modelo ajustado, repete-se todo o procedimento em cada uma, e o valor-$p$ é a fração das amostras sintéticas com distância KS maior que a observada. Por convenção, a lei de potência é rejeitada se $p < 0{,}1$. Para comparar a lei de potência com distribuições alternativas, como a lognormal ou a exponencial, usa-se a razão de verossimilhança normalizada de \citeonline{vuong1989}, cujo sinal indica o modelo favorecido e cujo valor-$p$ indica se a diferença é significativa. O pacote \texttt{powerlaw} implementa esse procedimento em Python \cite{alstott2014}.

\textbf{Estimador de Hill.} O estimador de \citeonline{hill1975} usa apenas as $N$ maiores observações $S_1 \geq S_2 \geq \dots \geq S_N$. Para a forma \textit{rank-size}, ele fornece diretamente
\begin{equation}
	\hat{\alpha}_N = \frac{1}{N - 1} \sum_{i=1}^{N-1} \ln \frac{S_i}{S_N},
	\qquad \hat{\zeta}_N = \frac{1}{\hat{\alpha}_N},
	\label{eq:hill}
\end{equation}
em que $\zeta = \gamma - 1$ é o expoente da distribuição complementar. O estimador de Hill é, essencialmente, o \gls{EMV} da \autoref{eq:mle} com $x_{\min} = S_N$, com $N-1$ no denominador em vez de $N$. A sua vantagem neste trabalho é a estabilidade: com $N$ fixo, o estimador sempre usa o mesmo número de observações, enquanto o $x_{\min}$ do método CSN pode saltar de forma descontínua entre amostras muito parecidas. A desvantagem é que o resultado depende de $N$, e por isso a dissertação o calcula para vários valores de $N$.

\section{Leis de escala urbana}
\label{sec:escala}

\citeonline{bettencourt2007} mostraram que muitos indicadores urbanos $Y$ variam com a população $X$ das cidades segundo leis de escala
\begin{equation}
	Y = c\, X^{\beta},
	\label{eq:escala}
\end{equation}
com expoentes que se agrupam em três regimes. Indicadores de infraestrutura, como a extensão de ruas e de redes elétricas, têm $\beta < 1$ (\textbf{sublinear}): cidades maiores precisam de menos infraestrutura por habitante. Necessidades individuais, como moradia e consumo de água, têm $\beta \approx 1$ (\textbf{linear}). Indicadores de interação social, como produção econômica, patentes e crimes, têm $\beta > 1$ (\textbf{superlinear}). \citeonline{bettencourt2013} propôs uma explicação para esses regimes a partir do equilíbrio entre os custos de deslocamento e os benefícios das interações numa rede urbana. Relações do mesmo tipo foram encontradas para emissões de CO$_2$ \cite{oliveira2014}, suicídios \cite{melo2014} e crimes contra o patrimônio em função da população flutuante \cite{caminha2017}.

Os expoentes de escala também dependem da definição de cidade. \citeonline{arcaute2015} mostraram que, no Reino Unido, o mesmo indicador pode parecer superlinear, linear ou sublinear conforme os limiares de densidade e de deslocamento usados para delimitar as cidades, e \citeonline{cottineau2017} chegaram a conclusão semelhante para a França. Esse resultado motiva, nesta dissertação, a escolha dos parâmetros do CCA por um critério externo, a lei de Zipf, e não pela qualidade do ajuste entre população e plástico.

\subsection{Ajuste e inferência}

A \autoref{eq:escala} é linear em escala logarítmica, $\log_{10} Y = \log_{10} c + \beta \log_{10} X$, e $\beta$ é tradicionalmente estimado por \gls{MQO} sobre os logaritmos \cite{bettencourt2007}. Com poucas observações e resíduos que não seguem uma distribuição normal, os intervalos de confiança analíticos do MQO não são confiáveis. Duas técnicas de reamostragem contornam esse problema. O \textit{bootstrap} de pares \cite{efron1979} sorteia, com reposição, $n$ pares $(X_i, Y_i)$ da amostra original, reajusta $\beta$ em cada amostra e usa os percentis 2,5 e 97,5 da distribuição resultante como \gls{IC} de 95\%. O teste de permutação \cite{good2005} avalia a hipótese nula $\beta = 0$: embaralhando os valores de $Y$ entre as observações, destrói-se qualquer associação com $X$, e o valor-$p$ é a fração das permutações com $|\beta|$ maior ou igual ao observado. Como medidas de associação, usam-se o coeficiente de correlação de Pearson $r$ entre os logaritmos, que mede a qualidade do ajuste linear em escala log-log, e o coeficiente de Spearman $\rho_s$, que mede a associação monotônica entre os postos e não depende da transformação logarítmica.
